\section{Related work}\label{sec:related}

\subsection{Three literatures, three questions}\label{sec:three_literatures}

The word ``reproducibility'' names different problems in different fields, and
the distinction determines what our measurement can and cannot speak to.

The broadest strand asks whether published findings are \emph{true}.
\citet{Ioannidis2005} argued on statistical grounds that most published research
findings are false, and the \citet{OpenScienceCollaboration2015} supplied direct
evidence by rerunning a hundred psychology experiments with fresh data. In
empirical finance the corresponding work is \citet{HarveyLiuZhu2016}, who
recalibrate significance thresholds for a literature that has tested hundreds of
candidate return predictors, and \citet{HouXueZhang2020}, who rebuild the
anomaly literature under uniform methodology and find much of it does not
survive. These studies concern the durability of findings under new data or
stricter statistics. They are not about whether the original code runs.

A second strand asks whether an artefact is \emph{documented} well enough that
reproduction is possible in principle. \citet{GundersenKjensmo2018} score
artificial-intelligence papers on documentation dimensions and find most
underspecified; the \citet{PineauReproducibility2021} NeurIPS programme turned
similar concerns into submission-time policy, introducing checklists and
artefact-availability expectations at a major venue. Instruments in this family,
including the disclosure rubric applied to machine-learning finance by
\citet{Khan2026Survey} and the audit-trail dimensions of
\citet{Khan2026MrAudit}, are scored by reading the paper and inspecting the
repository. They are cheap enough to apply at corpus scale, and they measure
disclosure.

The third question, and the one this paper takes up, is whether the released
artefact \emph{executes}. It sits between the other two: passing it is neither
necessary nor sufficient for a finding to be true, but failing it means no reader
can check the finding without reimplementing the work. The engineering literature
explains why this is hard even for competent authors.
\citet{SculleyTechnicalDebt2015} document how machine-learning systems accumulate
dependencies, configuration, and undeclared coupling faster than conventional
software, so a repository that ran on its author's machine drifts out of
executability without anyone touching it.

\subsection{Execution-based studies}\label{sec:exec_studies}

Two lines of work already run published code, and the present study is defined by
how it differs from each.

The first is at scale but outside finance. \citet{Trisovic2022} took more than
2,000 replication datasets from the Harvard Dataverse, comprising over 9,000 R
files, and executed them in a clean environment: 74\% of files failed on first
execution, falling to 56\% after automated code cleaning. That is the closest
methodological relative to what we do, and it establishes that large-scale
unassisted execution is both feasible and informative. It is confined to R and to
deposited replication packages rather than to repositories advertised in papers,
and it reports a single failure rate rather than locating where in the pipeline
the failure occurred.

The second is in finance but under very different conditions.
\citet{Perignon2024} ran the code of 168 research teams answering six
prespecified questions, and found that running that code on the same raw data
regenerates exactly the same results only 52\% of the time. This is the
definitive finance result on computational reproducibility, and it is a
\emph{best-case} measurement by construction: the code was written by
participating teams for the purpose, the raw data were supplied, and the study
was coordinated. Those are precisely the conditions an ordinary reader of a
published paper does not have.

Computer science has also institutionalised artefact evaluation at systems and
security venues, where committees run submitted artefacts and report substantial
attrition between what is advertised and what builds. That work is a human,
pass/fail judgement made by a reviewer who may correspond with the authors, which
makes it stronger evidence about an individual artefact and weaker evidence about
a distribution.

\subsection{What is new here}\label{sec:novelty}

The gap this paper fills is the unassisted, cold-start case in finance: not code
supplied for a coordinated study with data attached, but whatever a published
paper points a reader to, retrieved from that pointer, on a clean machine, with
no data supplied and no author contact. Read against \citet{Perignon2024}, the
question is not whether finance code reproduces under favourable conditions, but
how much of that 52\% survives the removal of the favourable conditions.

Three things are new. The ladder supplies an ordered scale where the literature
has used pass/fail or a single rate, so attrition becomes a curve with a
diagnosable shape rather than one number: our result is not simply that little
runs, but that most of the loss occurs at environment declaration and
construction rather than at link rot or at execution itself. The closed failure
vocabulary makes the \emph{reasons} a reportable table rather than a narrative.
And the harness is released as an instrument, which matters because the quantity
decays: a corpus scored today will score lower in two years as packages are
yanked and repositories are archived, and a study that cannot be re-run cannot
track that.
